% Part: normal-modal-logic
% Chapter: axioms-systems
% Section: logics-proofs

\documentclass[../../../include/open-logic-section]{subfiles}

\begin{document}

\olfileid{nml}{prf}{prf}

\olsection{\usetoken{P}{derivation} او مودالي نظامونه}

لومړے د عادي مودالي منطقونو دپاره !!a{derivation}
تعريفوو۔ په لنډ ډول، !!a{derivation} د !!{formula}s
داسې لړۍ ده چې هر !!{element} ئې يا د څو
\emph{بديهي اصلونو} څخه د يوه (د ځايناستۍ مثال)
وي، يا له مخکنيو !!{element}s څخه د څو استنتاجي
قاعدو له يوې راځي۔ د عادي مودالي منطقونو دپاره
د تاوتولوژيو\iftag{prvDiamond}{، \Ax{K}، او~\Dual{}}{ او~\Ax{K}}
ټول مثالونه بديهي اصلونه شمېرل کېږي۔ په دې سره
مودالي نظام~$\Log{K}$، يعنې تر ټولو کوچنے عادي
مودالي منطق، ترلاسه کېږي۔ خو د نورو نظامونو د
ترلاسه کولو دپاره نور بديهي اصلونه هم زياتولے شو۔
قاعدې تل مودوس پوننس~\MP{} او د لزوم ورکولو قاعده~\Nec دي۔

\begin{defn}
  که مودالي نظام~$\Log{K} !A_1 \dots !A_n$ او
  !!a{formula}~$!B$ راکړل شوي وي، وايو چې~$!B$
  په~$\Log{K} !A_1 \dots !A_n$ کښې \emph{!!{derivable}}
  دے، په نښو $\Log{K} !A_1 \dots !A_n \Proves !B$،
  هله او يوازې هله چې داسې !!{formula}s $!C_1$،
  \dots، $!C_k$ وي چې $!C_k = !B$، او هر~$!C_i$ يا
  د تاوتولوژۍ مثال وي، يا د $\Ax{K}$،\iftag{prvDiamond}{
  $\Dual$،}{} $!A_1$، \dots، $!A_n$ څخه د يوه مثال
  وي، يا د \MP{} يا~\Nec د قاعدې له مخې له مخکنيو
  !!{formula}s څخه راځي۔
\end{defn}

لاندې قضيه موږ ته اجازه راکوي چې د~$!B$ يو
$\Sigma$-!!{derivation} په ښودلو سره
$!B \in \Sigma$ ثابت کړو۔

\begin{prop}
  $\Log{K} !A_1 \dots !A_n = \Setabs{!B}{\Log{K} !A_1 \dots !A_n \Proves !B}$۔
\end{prop}

\begin{proof}
  د !!{derivation}s پر اوږدوالي استقرا کوو، څو وښيو چې
  $\Setabs{!B}{\Log{K} !A_1 \dots !A_n \Proves !B} \subseteq
  \Log{K} !A_1 \dots !A_n$۔

  که د~$!B$ د !!{derivation} اوږدوالے~$1$ وي، يوازې
  يو !!{formula} لري۔ دا !!{formula} د \MP{} يا \Nec
  په واسطه له مخکنيو فارمولونو څخه نۀ شي راتلے، نو
  بايد د تاوتولوژۍ مثال، د \Ax{K} مثال،\iftag{prvDiamond}{
  د $\Dual$ مثال،}{} يا د $!A_1$، \dots،~$!A_n$
  څخه د يوه مثال وي۔ خو $\Log{K}!A_1\dots!A_n$ هم
  دا لري، نو $!B \in \Log{K}!A_1 \dots !A_n$۔

  که د~$!B$ د !!{derivation} اوږدوالے~$> 1$ وي، نو
  $!B$ سربېره پر دې د \MP{} يا \Nec{} په واسطه له
  هغو !!{formula}s څخه هم ترلاسه کېدے شي چې د
  !!{derivation} وروستۍ کرښه نۀ وي۔ که $!B$ له~$!C$
  او~$!C \lif !B$ څخه (د \MP په واسطه) راځي، نو
  د استقرا د فرض له مخې $!C$ او $!C \lif !B \in
  \Log{K}!A_1 \dots !A_n$ دي۔ خو هر مودالي منطق
  د مودوس پوننس لاندې تړلے دے، نو
  $!B \in \Log{K}!A_1 \dots !A_n$۔ که
  $!B \equiv \Box !C$ د~$!C$ څخه د \Nec په واسطه
  راځي، نو د استقرا د فرض له مخې $!C
  \in \Log{K}!A_1 \dots !A_n$۔ خو هر عادي مودالي
  منطق د~$\Nec$ لاندې تړلے دے، نو
  $!B \in \Log{K}!A_1\dots!A_n$۔

  سرچپه شموليت داسې ښيو چې
  $\Sigma = \Setabs{!B}{\Log{K} !A_1 \dots !A_n \Proves !B }$
  يو عادي مودالي منطق دے چې د $!A_1$، \dots، $!A_n$
  ټول مثالونه لري؛ بيا دا کاروو چې
  $\Log{K} !A_1 \dots !A_n$ د تعريف له مخې تر ټولو
  کوچنے داسې منطق دے۔
  \begin{enumerate}
    \item هره تاوتولوژي~$!B$ د تاوتولوژۍ يو مثال دے، نو
      $\Log{K}!A_1\dots!A_n \Proves !B$، او $\Sigma$
      ټولې تاوتولوژۍ لري۔
    \item که $\Log{K}!A_1\dots!A_n \Proves !C$ او
      $\Log{K}!A_1\dots!A_n \Proves !C \lif !B$، نو
      $\Log{K}!A_1\dots!A_n \Proves !B$: د~$!C$
      !!{derivation} د~$!C \lif !B$ له اشتقاق
      سره يوځاے کړئ او د~$!B$ کرښه ورزياته کړئ۔
      وروستۍ کرښه د \MP{} له مخې توجيه کېږي۔ نو
      $\Sigma$ د مودوس پوننس لاندې تړلے دے۔
    \item که $!B$ يو !!a{derivation} ولري، نو د~$!B$
      هر د ځايناستۍ مثال هم !!a{derivation} لري:
      ځايناستي د !!{derivation} پر هر !!{formula}
      تطبيق کړئ۔ (تمرین: د !!{derivation}s پر اوږدوالي
      په استقرا ثابت کړئ چې پايله لا هم سم
      !!{derivation} دے۔) نو $\Sigma$ د يو شان
      ځايناستۍ لاندې تړلے دے۔ (اوس مو ثابته کړه
      چې $\Sigma$ د مودالي منطق ټول شرطونه پوره کوي۔)
    \item $\Log{K}!A_1\dots!A_n \Proves \Ax{K}$ لرو، نو
      $K \in \Sigma$۔
    \tagitem{prvDiamond}{همدارنګه
      $\Log{K}!A_1\dots!A_n \Proves \Dual$ لرو، نو
      $\Dual \in \Sigma$۔}{}
    \item که $\Log{K}!A_1\dots!A_n \Proves !C$، نو
      زياته شوې کرښه~$\Box !C$ د~\Nec له مخې
      توجيه کېږي۔ په پايله کښې $\Sigma$ د~\Nec
      لاندې تړلے دے۔ نو $\Sigma$ عادي دے۔
  \end{enumerate}
\end{proof}

\end{document}
